\documentclass[10pt]{beamer}

% --- Preamble and Packages ---
\usepackage[utf8]{inputenc}
\usepackage{graphicx}
\usepackage{amsmath}
\usepackage{amssymb}
\usepackage{hyperref}
\usepackage{booktabs}
\usepackage{tikz}
\usetikzlibrary{shapes.geometric, arrows, positioning, chains, fit, calc}
\usepackage{epstopdf}

\usetheme{CambridgeUS}
\usecolortheme{dolphin}

% --- Custom Colors ---
\definecolor{myNewColorA}{RGB}{90, 80, 140}
\setbeamercolor*{palette primary}{bg=myNewColorA}
\setbeamercolor*{titlelike}{fg=myNewColorA}
\setbeamercolor*{title}{bg=myNewColorA, fg=white}
\setbeamercolor*{item}{fg=myNewColorA}
\setbeamercolor*{caption name}{fg=myNewColorA}
\usefonttheme{professionalfonts}

% --- Font Sizes ---
\setbeamerfont{title}{size=\large}
\setbeamerfont{subtitle}{size=\small}
\setbeamerfont{author}{size=\small}
\setbeamerfont{date}{size=\small}
\setbeamerfont{institute}{size=\footnotesize}

% --- Title Info ---
\title[CFO Estimation under Impulsive Noise]{CFO Estimation for OFDM-Based Wireless Communication Systems Under Impulsive Noise}
\author[Shivansh Rai et al.]{Shivansh Rai\inst{1}, Sandesh Jain\inst{2}, Mahesh Shamrao Chaudhari\inst{3}, \and Poornima Singh Thakur\inst{1}}
\institute[]{
    \inst{1} \textit{Dept. of Information Technology, ABV-IIITM Gwalior} \\
    \inst{2} \textit{Dept. of Electrical and Electronics Engineering, ABV-IIITM Gwalior} \\
    \inst{3} \textit{Dept. of Electronics and Communication Engineering, TIET Patiala}
}
\date{March 13 , 2026}

% =====================================================================
\begin{document}
% =====================================================================

\frame{\titlepage}

\begin{frame}
    \frametitle{Outline}
    \tableofcontents
\end{frame}

% =====================================================================
\section{Introduction}
% =====================================================================

\begin{frame}
\frametitle{Motivation and Problem Statement}

\begin{block}{The Impact of CFO}
Modern wireless standards — \textbf{4G LTE, 5G NR, Wi-Fi (802.11ax)} — all rely on
orthogonal frequency division multiplexing (OFDM) for its high spectral efficiency
and resilience to multipath fading.
However, OFDM is highly sensitive to synchronization errors:
carrier frequency offset (CFO), caused by oscillator mismatch or Doppler shifts,
destroys subcarrier orthogonality and leads to severe inter-carrier interference (ICI),
significantly degrading bit error rate (BER) performance.
\end{block}

\vspace{0.3cm}

\begin{alertblock}{Limitations of Existing DL Estimators}
Most deep learning (DL) approaches for CFO estimation rely on the conventional
minimum mean squared error (MMSE) criterion.
\begin{itemize}
    \item MMSE minimizes average squared error, relying on second-order statistics.
    \item Under non-Gaussian or impulsive noise, squared-error terms \textbf{amplify} high-amplitude outliers.
    \item This causes \textbf{gradient explosion} during training, hindering convergence.
\end{itemize}
\end{alertblock}

\end{frame}

% =====================================================================
\section{System and Noise Model}
% =====================================================================

\begin{frame}
\frametitle{System Model Overview}
\begin{figure}
    \centering
    \includegraphics[width=\textwidth, height=0.72\textheight, keepaspectratio]{system_model.eps}
    \caption{Signal propagation: transmitter, channel, CFO, impulsive noise, and CAD-Huber estimator.}
\end{figure}
\end{frame}

% \begin{frame}
% \frametitle{OFDM Signal Model}

% \begin{block}{Transmitted Signal}
% Time-domain signal $x[n]$ after IFFT and CP insertion:
% \[
% x[n]=\frac{1}{\sqrt{N}}\sum_{k=0}^{N-1}X[k]\,e^{j\frac{2\pi kn}{N}}, \qquad -N_{cp}\le n<N
% \]
% \end{block}

% \vspace{0.3cm}

% \begin{block}{Received Signal}
% After a frequency-selective channel with $L$ taps and normalized CFO $\epsilon$:
% \[
% r[n]=e^{j\frac{2\pi\epsilon n}{N}}\sum_{l=0}^{L-1}h[l]\,x[n-l]+w[n]
% \]
% \end{block}

% \end{frame}

\begin{frame}
\frametitle{Impulsive Noise Model}

\begin{alertblock}{Symmetric $\alpha$-Stable ($S\alpha S$) Distribution}
We model $w[n]$ using the $S\alpha S$ distribution to capture impulsive industrial/urban channels.
Its characteristic function is:
\[
\varphi(t)=\exp\!\bigl(-\gamma\,|t|^{\alpha}\bigr)
\]
\end{alertblock}

\vspace{0.3cm}

\begin{columns}[T]
    \column{0.5\textwidth}
    \begin{block}{Parameters}
    \begin{itemize}
        \item $\alpha \in (0,2]$: \textit{characteristic exponent} — controls outlier severity.
        \item $\gamma > 0$: \textit{dispersion parameter}.
    \end{itemize}
    \end{block}

    \column{0.5\textwidth}
    \begin{block}{Key Implication}
    For $\alpha < 2$, variance is \textbf{infinite}.
    Conventional $L_2$-norm estimators diverge — necessitating a robust alternative.
    \end{block}
\end{columns}

\end{frame}

% --- Architecture Part 1: Diagram ---
\begin{frame}
\frametitle{CAD-HuberNet Architecture}
\begin{figure}
    \centering
    \includegraphics[width=\textwidth, height=0.72\textheight, keepaspectratio]{architecture.eps}
    \caption{Detailed view of the dual-stream CNN-Attention-DNN (CAD-Huber) model.}
\end{figure}
\end{frame}
% =====================================================================
\section{Proposed Architecture: CAD-HuberNet}
% =====================================================================

\begin{frame}
\frametitle{Multi-Modal Input Feature Representation}

\begin{exampleblock}{Dual-Stream Inputs: $\mathbf{X} = [X^{(u)},\, X^{(l)}]$}

\vspace{0.15cm}
\begin{columns}[T]
    \column{0.5\textwidth}
    \textbf{Upper Stream} (Squared Spectrum)\\[4pt]
    Amplifies cyclic frequencies linked to carrier offset:
    \[
    X^{(u)}=\bigl|\mathcal{F}\!\left(|r[n]|^{2}\right)\bigr|\in\mathbb{R}^{S\times1}
    \]

    \column{0.5\textwidth}
    \textbf{Lower Stream} (Magnitude + Phase)\\[4pt]
    Preserves phase information to resolve frequency ambiguities:
    \[
    X^{(l)} = \bigl[|\mathcal{F}(r[n])|,\; \angle\mathcal{F}(r[n])\bigr] \in\mathbb{R}^{S\times2}
    \]
\end{columns}
\end{exampleblock}

\vspace{0.2cm}
Raw time-domain signals lack direct spectral features; the dual-stream design ensures both
power-spectral and phase information are fed to the network simultaneously.

\end{frame}



% --- Architecture Part 2: Component Details ---
\begin{frame}
\frametitle{CAD-HuberNet: Key Components}

\begin{block}{1. Feature Extraction — Triple Residual Stacks (TRS)}
Prevents vanishing gradients via skip connections. For the $l$-th layer:
\[
Z_{l}=\sigma(W_{l}*Z_{l-1}+b_{l})+Z_{l-1}
\]
\end{block}

\begin{block}{2. Attention Mechanism}
Prioritizes spectral features from clean subcarriers; suppresses impulsive-spike-corrupted ones:
\[
e_{i}=v^{T}\tanh(W_{a}h_{i}+b_{a}), \qquad
\alpha_{i}=\frac{\exp(e_{i})}{\displaystyle\sum_{k}\exp(e_{k})}
\]
\end{block}

\begin{block}{3. DNN Regression Head}
Attention-weighted context vectors from both streams are concatenated and passed through
fully-connected layers to produce the scalar CFO estimate $\hat{\epsilon}$.
\end{block}

\end{frame}

% =====================================================================
\section{Robust Optimization}
% =====================================================================

\begin{frame}
\frametitle{Robust Optimization via Huber Loss}

\begin{block}{The Huber Penalty Function}
To mitigate exploding gradients from heavy-tailed $S\alpha S$ noise, the network is trained
with the Huber loss:
\[
\rho_{\delta}(e)=
\begin{cases}
\dfrac{1}{2}e^{2} & |e|\le\delta,\\[8pt]
\delta|e|-\dfrac{1}{2}\delta^{2} & |e|>\delta.
\end{cases}
\]
\end{block}

\vspace{0.2cm}

\begin{alertblock}{Gradient Saturation — Preventing Explosion}
The gradient $\psi(e)$ clips large residuals to a fixed magnitude $\delta$:
\[
\psi(e) =
\begin{cases}
e & |e|\le\delta,\\[4pt]
\delta\cdot\operatorname{sgn}(e) & |e|>\delta.
\end{cases}
\]
Massive outliers are therefore \textbf{bounded} during backpropagation, ensuring stable convergence.
\end{alertblock}

\end{frame}

% =====================================================================
\section{Simulation Results}
% =====================================================================

\begin{frame}
\frametitle{Simulation Setup}

\begin{columns}[T]
    \column{0.5\textwidth}
    \begin{block}{OFDM System Parameters}
    \begin{itemize}
        \item Subcarriers: $N = 128$
        \item Cyclic prefix: $N_{cp} = 16$
        \item Modulation: 64-QAM
        \item Channel taps: $L = 3$ (freq.-selective)
        \item Normalized CFO: $\epsilon \sim \mathcal{U}(-0.5,\,0.5)$
    \end{itemize}
    \end{block}

    \column{0.5\textwidth}
    \begin{block}{Training Configuration}
    \begin{itemize}
        \item Noise: $S\alpha S$, $\alpha=1.4$, $\gamma=0.1$
        \item Dataset: 20{,}000 OFDM symbols
        \item Split: 70\% train / 30\% validation
        \item Optimizer: Adam, LR $= 10^{-3}$
        \item Batch size: 1000; max epochs: 2000
    \end{itemize}
    \end{block}
\end{columns}

\vspace{0.3cm}
\begin{block}{Baseline}
All results compare \textbf{CAD-HuberNet} against the existing \textbf{CAD-MMSE} model
\cite{chaudhari2023_cad}. Performance metric: MSE (dB scale, $10\log_{10}(\text{MSE})$).
\end{block}

\end{frame}

\begin{frame}
\frametitle{Results: Convergence and SNR Performance}

\begin{columns}[T]

\column{0.5\textwidth}
\begin{figure}
\centering
\includegraphics[width=\linewidth, height=0.48\textheight, keepaspectratio]{mse_vs_epoch.eps}
\caption{Training Convergence (MSE vs. Epoch)}
\end{figure}
\footnotesize
MMSE loss oscillates due to impulsive spikes amplifying the squared-error gradient.
Huber loss yields \textbf{stable, monotone} descent.

\column{0.5\textwidth}
\begin{figure}
\centering
\includegraphics[width=\linewidth, height=0.48\textheight, keepaspectratio]{steady_mse_vs_snr.eps}
\caption{Steady-State MSE vs. SNR}
\end{figure}
\footnotesize
At 15 dB SNR, CAD-HuberNet achieves \textbf{$-19.21$ dB MSE} vs.\ $-9.42$ dB for CAD-MMSE
— a gain of nearly \textbf{10 dB}.

\end{columns}

\end{frame}

\begin{frame}
\frametitle{Results: Robustness to Impulsiveness ($\alpha$)}

\begin{figure}
\centering
\includegraphics[width=0.65\textwidth, height=0.5\textheight, keepaspectratio]{mse_vs_alpha.eps}
\caption{Steady-state MSE vs.\ $\alpha$ at SNR = 20 dB}
\end{figure}

\vspace{0.1cm}
\begin{itemize}
    \item In the \textbf{highly impulsive} regime ($\alpha=1.3$), CAD-HuberNet achieves
          MSE $\approx -1.38$ dB vs.\ CAD-MMSE at MSE $\approx 0$ dB.
    \item As $\alpha \to 2.0$ (Gaussian limit), both models converge to comparable performance,
          confirming Huber loss incurs \textbf{no penalty} in benign conditions.
\end{itemize}

\end{frame}

% =====================================================================
\section{Conclusion}
% =====================================================================

\begin{frame}
\frametitle{Conclusion}

\begin{block}{Summary of Contributions}
\begin{itemize}
    \item Proposed a \textbf{robust blind CFO estimation} framework for heavy-tailed
          non-Gaussian (impulsive) noise environments.
    \item Designed a \textbf{dual-stream CNN-Attention-DNN (CAD)} architecture that
          jointly exploits power-spectral and phase features.
    \item Replaced MMSE with \textbf{Huber loss}, actively bounding outlier gradients
          during backpropagation.
    \item Demonstrated \textbf{superior steady-state accuracy and convergence stability}
          over standard DL baselines, especially for $\alpha < 2$.
\end{itemize}
\end{block}

\end{frame}

% --- THANK YOU SLIDE ---
\begin{frame}[plain]
    \centering
    \vfill
    {\Huge \bfseries Thank You}
    \vspace{1.2em}
    \large Questions?
    \vfill
    \normalsize
    \textbf{Shivansh Rai}\\[4pt]
    \footnotesize Department of Information Technology\\
    ABV-IIITM Gwalior
    \vfill
\end{frame}

\end{document}
